% Options for packages loaded elsewhere
\PassOptionsToPackage{unicode}{hyperref}
\PassOptionsToPackage{hyphens}{url}
%
\documentclass[
  11pt,
]{article}
\usepackage{amsmath,amssymb}
\usepackage{iftex}
\ifPDFTeX
  \usepackage[T1]{fontenc}
  \usepackage[utf8]{inputenc}
  \usepackage{textcomp} % provide euro and other symbols
\else % if luatex or xetex
  \usepackage{unicode-math} % this also loads fontspec
  \defaultfontfeatures{Scale=MatchLowercase}
  \defaultfontfeatures[\rmfamily]{Ligatures=TeX,Scale=1}
\fi
\usepackage{lmodern}
\ifPDFTeX\else
  % xetex/luatex font selection
\fi
% Use upquote if available, for straight quotes in verbatim environments
\IfFileExists{upquote.sty}{\usepackage{upquote}}{}
\IfFileExists{microtype.sty}{% use microtype if available
  \usepackage[]{microtype}
  \UseMicrotypeSet[protrusion]{basicmath} % disable protrusion for tt fonts
}{}
\makeatletter
\@ifundefined{KOMAClassName}{% if non-KOMA class
  \IfFileExists{parskip.sty}{%
    \usepackage{parskip}
  }{% else
    \setlength{\parindent}{0pt}
    \setlength{\parskip}{6pt plus 2pt minus 1pt}}
}{% if KOMA class
  \KOMAoptions{parskip=half}}
\makeatother
\usepackage{xcolor}
\usepackage[margin=1in]{geometry}
\usepackage{color}
\usepackage{fancyvrb}
\newcommand{\VerbBar}{|}
\newcommand{\VERB}{\Verb[commandchars=\\\{\}]}
\DefineVerbatimEnvironment{Highlighting}{Verbatim}{commandchars=\\\{\}}
% Add ',fontsize=\small' for more characters per line
\newenvironment{Shaded}{}{}
\newcommand{\AlertTok}[1]{\textcolor[rgb]{1.00,0.00,0.00}{\textbf{#1}}}
\newcommand{\AnnotationTok}[1]{\textcolor[rgb]{0.38,0.63,0.69}{\textbf{\textit{#1}}}}
\newcommand{\AttributeTok}[1]{\textcolor[rgb]{0.49,0.56,0.16}{#1}}
\newcommand{\BaseNTok}[1]{\textcolor[rgb]{0.25,0.63,0.44}{#1}}
\newcommand{\BuiltInTok}[1]{\textcolor[rgb]{0.00,0.50,0.00}{#1}}
\newcommand{\CharTok}[1]{\textcolor[rgb]{0.25,0.44,0.63}{#1}}
\newcommand{\CommentTok}[1]{\textcolor[rgb]{0.38,0.63,0.69}{\textit{#1}}}
\newcommand{\CommentVarTok}[1]{\textcolor[rgb]{0.38,0.63,0.69}{\textbf{\textit{#1}}}}
\newcommand{\ConstantTok}[1]{\textcolor[rgb]{0.53,0.00,0.00}{#1}}
\newcommand{\ControlFlowTok}[1]{\textcolor[rgb]{0.00,0.44,0.13}{\textbf{#1}}}
\newcommand{\DataTypeTok}[1]{\textcolor[rgb]{0.56,0.13,0.00}{#1}}
\newcommand{\DecValTok}[1]{\textcolor[rgb]{0.25,0.63,0.44}{#1}}
\newcommand{\DocumentationTok}[1]{\textcolor[rgb]{0.73,0.13,0.13}{\textit{#1}}}
\newcommand{\ErrorTok}[1]{\textcolor[rgb]{1.00,0.00,0.00}{\textbf{#1}}}
\newcommand{\ExtensionTok}[1]{#1}
\newcommand{\FloatTok}[1]{\textcolor[rgb]{0.25,0.63,0.44}{#1}}
\newcommand{\FunctionTok}[1]{\textcolor[rgb]{0.02,0.16,0.49}{#1}}
\newcommand{\ImportTok}[1]{\textcolor[rgb]{0.00,0.50,0.00}{\textbf{#1}}}
\newcommand{\InformationTok}[1]{\textcolor[rgb]{0.38,0.63,0.69}{\textbf{\textit{#1}}}}
\newcommand{\KeywordTok}[1]{\textcolor[rgb]{0.00,0.44,0.13}{\textbf{#1}}}
\newcommand{\NormalTok}[1]{#1}
\newcommand{\OperatorTok}[1]{\textcolor[rgb]{0.40,0.40,0.40}{#1}}
\newcommand{\OtherTok}[1]{\textcolor[rgb]{0.00,0.44,0.13}{#1}}
\newcommand{\PreprocessorTok}[1]{\textcolor[rgb]{0.74,0.48,0.00}{#1}}
\newcommand{\RegionMarkerTok}[1]{#1}
\newcommand{\SpecialCharTok}[1]{\textcolor[rgb]{0.25,0.44,0.63}{#1}}
\newcommand{\SpecialStringTok}[1]{\textcolor[rgb]{0.73,0.40,0.53}{#1}}
\newcommand{\StringTok}[1]{\textcolor[rgb]{0.25,0.44,0.63}{#1}}
\newcommand{\VariableTok}[1]{\textcolor[rgb]{0.10,0.09,0.49}{#1}}
\newcommand{\VerbatimStringTok}[1]{\textcolor[rgb]{0.25,0.44,0.63}{#1}}
\newcommand{\WarningTok}[1]{\textcolor[rgb]{0.38,0.63,0.69}{\textbf{\textit{#1}}}}
\setlength{\emergencystretch}{3em} % prevent overfull lines
\providecommand{\tightlist}{%
  \setlength{\itemsep}{0pt}\setlength{\parskip}{0pt}}
\setcounter{secnumdepth}{-\maxdimen} % remove section numbering
\ifLuaTeX
  \usepackage{selnolig}  % disable illegal ligatures
\fi
\usepackage{bookmark}
\IfFileExists{xurl.sty}{\usepackage{xurl}}{} % add URL line breaks if available
\urlstyle{same}
\hypersetup{
  pdftitle={Curve recognition and actual local analytic images},
  pdfauthor={P06.2 reviewed partial result},
  hidelinks,
  pdfcreator={LaTeX via pandoc}}

\title{Curve recognition and actual local analytic images}
\author{P06.2 reviewed partial result}
\date{3 October 2026}

\begin{document}
\maketitle

\subsection{Status, conventions, and exact
statements}\label{status-conventions-and-exact-statements}

This paper gives independently reviewed partial results for P06.2. The
exact proof, its complete supporting argument and authentic primary
source bodies received the source-aware PASS accompanying this paper.
The original problem remains open. No formal verification or priority
certification is claimed; earlier failed and inconclusive records are
preserved.

All spaces and morphisms are Berkovich analytic spaces over the field
specified in each statement. An ordinary open is open in the underlying
topology. An analytic neighbourhood contains an ordinary neighbourhood;
an analytic domain alone need not do so. A curve is separated and pure
of analytic dimension one. Put \(D=d'd''\). All forms and currents in an
equation below are \textbf{ordinary} Chambert--Loir--Ducros (CLD) forms
and currents, defined on ordinary opens. For a pure space \(Z\),
\([u]=u\delta_Z\) uses canonical integration with the actual generic
artinian lengths of its structure, including nilpotents. In dimension
one, \[
                 D[u](\psi)=\int_Z uD\psi,
                 \qquad \psi\in A_c^{0,0}(Z).
\] The equation is tested against every compact ordinary smooth scalar
test. Smooth top forms define signed Radon measures, so these integrals
with a continuous multiplier are well defined. This is not an a priori
assertion that the differentiated current \(D[u]\) is a measure.

Here ``harmonic'\,' on an arbitrary curve means Gubler--Rabinoff (GR)
Definition 3.13: both signs are subharmonic in GR Definition 3.4, which
includes analytic field extensions. On a smooth strict curve over a
nontrivially valued field this agrees with Thuillier harmonicity.
Classical pluriharmonicity means that both signs are classically psh in
GR Definition 4.1, with its \textbf{every analytic extension and every
analytic curve map} quantifiers. Gubler--Jell--Rabinoff (GJR)
harmonicity is the real-PL residual numerical-triviality notion; it is
not ordinary CLD smoothness. Real PL is used in the precise GR 2.16/GJR
convention: G-locally on strict analytic domains the function is a
finite real linear combination of logarithms of analytic units. The
local real-PL conclusions below refer to that convention, not to one
ordinary smooth moment expression.

\textbf{Theorem A (intrinsic curve recognition).} Let \(K\) be
algebraically closed, complete, and nontrivially valued. Let \(Z\) be a
good separated strict boundaryless pure analytic curve. If
\(u:Z\to\mathbb R\) is continuous and \(D[u]=0\), then \(u\) is
GR-harmonic, classically pluriharmonic, locally real PL, and GJR
harmonic. If \(Z\) is smooth, \(u\) is Thuillier harmonic. There is no
reducedness, normality, properness, or global algebraic-model
hypothesis. The theorem also applies on every ordinary open of \(Z\).

\textbf{Proposition B (local images, arbitrary field).} Let \(K\) be
complete and nontrivially valued, \(X\) good separated strict
boundaryless, and \(f:Y\to X\) a \(K\)-analytic morphism from a good
separated smooth strict boundaryless pure curve. At each analytic point
\(y\in Y\), either:

\begin{enumerate}
\def\labelenumi{\arabic{enumi}.}
\tightlist
\item
  on an ordinary neighbourhood of \(y\), \(f\) factors analytically
  through \(\mathcal M(E)\) for a finite field extension \(E/K\),
  mapping to a single rigid point of \(X\); or
\item
  there are strict affinoid neighbourhoods \(V\) of \(y\) and \(U\) of
  \(f(y)\), with both points in their ordinary interiors, such that
  \(V\) is connected and normal and \(V\to U\) is finite. Its reduced
  kernel image \(C\hookrightarrow U\) is irreducible and pure
  one-dimensional, \(V\to C\) is finite surjective, and it admits a
  finite analytic lift to the finite normalization of \(C\).
\end{enumerate}

\textbf{Theorem C (ground-field ambient probes).} Let \(K\) be
algebraically closed, complete, and nontrivially valued. Let \(X\) be
good separated strict boundaryless and pure of dimension \(n\geq1\). If
\(h:X\to\mathbb R\) is continuous with \(D[h]=0\), then, for every
\(K\)-analytic map \(f:Y\to X\) as in Proposition B, \(h\circ f\) is
Thuillier harmonic on \(Y\).

Theorem C is fieldwise. It does \textbf{not} say that \(h\) is
classically pluriharmonic on \(X\), or that the equation ascends to
\(X_L\). Even when \(K\) is algebraically closed, analytic extensions
\(L/K\) can be transcendental. The universal conclusion of Theorem A is
justified instead by the already-recognized curve harmonicity and the
precise GR curve functoriality theorem.

\subsection{Exact source inputs and
applicability}\label{exact-source-inputs-and-applicability}

The primary PDFs and their provenance are copied in \texttt{sources/}.
Fresh local PDF extractions are in \texttt{evidence/}; numbered selected
passages are in \texttt{sources/primary-passages.txt}. The source and
artifact manifests record origins, byte sizes, and SHA-256 values.
Cached provenance identifies the source URL; this worker claims no new
network retrieval. The inspected versions, not hypothetical later
versions, determine all numbering below.

\subsubsection{S1. Supplied supporting theorem
n100}\label{s1.-supplied-supporting-theorem-n100}

The imported note n100 supplies \emph{Intrinsic zero-Hessian restriction
to closed analytic curves}, frozen argument SHA-256
\texttt{be42c6a0e7669dbf1c6fe33ad46ee1f0c020a66efb01a67abff6ee3845e6f836}.
The local copy equals that hash. Its complete proof and earlier exact-scope
review remain preserved as \texttt{sources/closed-curves-argument.md} and
\texttt{sources/closed-curves-review.json}. The present review argument
also embeds and independently reconstructs the entire supporting proof;
the expanded argument and full review receipt accompany this paper.
The precise results used
are:

\begin{itemize}
\tightlist
\item
  Over any complete nontrivially valued \(K\), if \(M\) is good
  separated strict boundaryless pure of dimension \(n\geq1\),
  \(\Omega\subset M\) is ordinary open, and \(h\) is continuous with
  \(D(h\delta_\Omega)=0\), then for every closed pure analytic curve
  \(C\hookrightarrow M\), \(D(h|_{C\cap\Omega}\delta_{C\cap\Omega})=0\)
  intrinsically. This is also local for a closed curve in an ordinary
  open. Both structures may be nonreduced and all actual generic lengths
  are retained.
\item
  On a good separated strict boundaryless reduced pure curve, a
  continuous zero-current equation extends across a locally finite set
  of rigid punctures.
\item
  For the finite normalization \(\nu:\widetilde C\to C\) of such a
  reduced curve, \(D[g]=0\) implies \(D[g\circ\nu]=0\). The proof uses
  the actual isomorphism locus and then rigid-puncture removal. It does
  not invert a finite direct image or give arbitrary finite-cover
  current pullback.
\end{itemize}

We use these results without reproving them and declare n100 as native
support wherever they are used. I assessed their proof bodies: their
cutting argument passes through reduced \(S_1\) component cycles,
positive weak limits multiplied by a continuous function, and
closed-embedding test lifts; it does not declare an arbitrary
embedded-prime quotient \(S_1\). The normalization exception is removed
by a continuous puncture argument. In Theorem A they are used on
reduction and normalization only; in Theorem C they are used on an
actual closed local image inside an ordinary boundaryless open. None of
their results gives recognition or general scalar-current ascent.

\subsubsection{S2. Affinoid neighbourhoods, algebra, dimensions, and
smoothness}\label{s2.-affinoid-neighbourhoods-algebra-dimensions-and-smoothness}

CLD means A. Chambert-Loir and A. Ducros, arXiv:1204.6277v2, 25 July
2025. The exact consequences needed here are these.

\begin{itemize}
\tightlist
\item
  CLD 11.1.2: an affinoid algebra is noetherian, its ideals are closed,
  \(\mathcal M(A)\to\operatorname{Spec}(A)\) is surjective, and the
  spectral seminorm vanishes precisely on the nilpotents. On a reduced
  affinoid algebra it is an equivalent Banach norm. A strict affinoid
  algebra has a bounded quotient presentation from a Tate algebra with
  unit radii.
\item
  CLD 11.1.5--7: a good analytic space has an affinoid-neighbourhood
  basis, and ordinary opens are analytic domains. CLD 11.7.7 explicitly
  adds that every point of a good strict space over a nontrivially
  valued field has a \textbf{strict affinoid neighbourhood}. Applying it
  inside a prescribed ordinary open gives such a neighbourhood there.
\item
  CLD 11.2.1--3: a closed analytic subspace of \(\mathcal M(A)\) is
  \(\mathcal M(A/I)\) for an ideal \(I\); on an affinoid the Zariski
  closed sets are such zero sets. The reduction has the same carrier.
  Analytic domains preserve local analytic dimension, and the dimension
  of a strict affinoid equals the Krull dimension of its algebra. CLD
  11.2.5: affinoids have finitely many irreducible components; an
  irreducible analytic space is pure, and its proper closed analytic
  subsets have smaller dimension.
\item
  CLD 11.3.1: the stated local-ring regularity conditions persist on
  analytic domains; the normalization is finite surjective, and locally
  is the integral closure of the reduced algebra in its total quotient
  ring. CLD 11.4.8: quasi-smoothness is geometric regularity, hence
  implies regularity and normality here. CLD 11.4.9: on good spaces
  smoothness means quasi-smoothness together with relative interior.
  Thus a boundaryless rig-smooth curve is smooth.
\end{itemize}

We also use the strict affinoid Nullstellensatz, in its exact form from
Bosch--Güntzer--Remmert (BGR), \emph{Non-Archimedean Analysis}, 6.1.2/1:
over a complete nontrivially valued field, the residue field at a
maximal ideal of a strict affinoid algebra is finite over the ground
field. Thus every nonempty Zariski closed subset of a strict affinoid
contains a rigid point. This follows by choosing a maximal ideal of its
nonzero quotient and using the finite-extension valuation. A
zero-dimensional strict affinoid algebra is finite over \(K\): it is
noetherian artinian, its semisimple quotient is a finite product of the
preceding finite fields, and its nilpotent radical has a finite
filtration with finite-dimensional successive quotients. A fresh BGR PDF
is not in the supplied cache; this standard primary theorem is
explicitly stated and cited, not claimed to have been verified by a
PDF-byte check. Its strict/nontrivial hypotheses are checked at every
application.

The other valued-field fact is BGR 3.2.4/2: a finite algebraic extension
of a complete non-Archimedean field has a unique extending absolute
value and is complete. It includes inseparable extensions. In the
application to a finite field inside a reduced affinoid, uniqueness also
follows by equivalence of finite-dimensional ultrametric norms: for two
extending norms, \(|z|_1\leq C|z|_2\); applying this to \(z^m\) and
taking \(m\)th roots gives \(|z|_1\leq|z|_2\), and then the reverse
inequality. The usual finite-dimensional completeness and
norm-equivalence proof is by induction, using that a complete
codimension-one subspace is closed and has positive distance from any
fixed vector outside it.

\subsubsection{S3. Boundary and local finiteness, with different
quantifiers}\label{s3.-boundary-and-local-finiteness-with-different-quantifiers}

CLD 11.4.1 states that for an analytic-domain inclusion the relative
interior is the ordinary topological interior. An algebraic
analytification is boundaryless. For \(f:Y\to X\) with \textbf{separated
target} \(X\), formula 11.4.1.3 states exactly \[
 \operatorname{Int}(Y/X)\cap f^{-1}(\operatorname{Int}(X))
       =\operatorname{Int}(Y),\qquad
 \partial Y=\partial(Y/X)\cup f^{-1}(\partial X).
\] Boundaryless maps are stable under composition and analytic base
change. The equality is also proved in the author-hosted Conrad--Temkin,
\emph{Descent for non-Archimedean analytic spaces}, arXiv:1912.06230v2,
Lemma 2.9: for \(\psi:A\to B\) and \(\phi:B\to C\), the inclusion
\(\operatorname{Int}(A/B)\cap\psi^{-1}\operatorname{Int}(B/C)
\subset\operatorname{Int}(A/C)\) is equality when \(\phi\) is locally
separated. We inspected its actual proof. Only this boundary equality,
not a current-ascent or general descent result from that paper, is used.

CLD 11.4.4 states that finite maps are boundaryless, finite maps to an
affinoid have affinoid inverse images, and finite analytic maps are
equivalent to finite Banach algebra maps on affinoids. Finite Banach
\(A\)-algebras and finite ring \(A\)-algebras are equivalent categories;
the finite-module quotient topology determines the Banach topology. Its
footnote explicitly adds: if \(|K^\times|\ne\{1\}\) and \(A\) is strict,
a Banach \(A\)-algebra finite as a ring algebra is finite Banach. A map
is \textbf{finite at \(y\)} precisely when it is finite from an analytic
neighbourhood \(V\) of \(y\) to an analytic neighbourhood \(U\) of its
image.

CLD 11.4.5 states that \(\dim_y f=\dim_y(Y_{f(y)})\), with the actual
fibre over \(\mathcal H(f(y))\), is Zariski upper semicontinuous (citing
Ducros, \emph{Variation de la dimension relative}, Theorem 4.9), and
that \(f\) is finite at \(y\) if and only if \(\dim_y f=0\) and
\(y\in\operatorname{Int}(Y/X)\). This is the local criterion used below.
CLD 11.6.6 concerns finiteness over a target neighbourhood of an
\textbf{entire fibre} and has a compactness hypothesis. We neither use
that different criterion nor assert compactness of an arbitrary original
fibre.

\subsubsection{S4. The checked affinoid embedding
theorem}\label{s4.-the-checked-affinoid-embedding-theorem}

M. van der Put, \emph{The class group of a one-dimensional affinoid
space}, Ann. Inst. Fourier 30(4) (1980), 155--164, Theorem 1.1, printed
page 156:

\begin{quote}
Every 1-dimensional, normal, connected affinoid space
\(Y=\mathrm{Sp}(A)\) is an affinoid subspace of a non-singular complete
curve.
\end{quote}

The field assumption on printed page 155 is that \(k\) is
\textbf{complete and algebraically closed}, with a non-Archimedean
valuation. Its affinoid convention is an actual quotient of
\(k\langle T_1,\ldots,T_m\rangle\); these are strict affinoids. The
proof defines normal/connected by integral closedness and absence of
nontrivial idempotents. It glues finitely many affinoid pieces to \(Y\)
and explicitly concludes that \(Y\) is an affinoid domain of the
resulting \textbf{algebraic} complete nonsingular curve. Over our
algebraically closed field that curve is smooth and proper. We use this
theorem only over the same algebraically closed ground field.

For clarity about analytic categories, the exact Berkovich-domain
formulation is also explicitly used in Thuillier, HAL tel-00010990v1,
the proof of Theorem 2.3.8, printed page 53: a pure one-dimensional
rig-smooth strict affinoid, after a finite separable extension
\(k_1/k\), is isomorphic to an affinoid domain of a smooth proper
algebraic curve over \(k_1\), with a citation to van der Put's theorem.
This is an actual assertion in that proof, not an inference from its
bibliography. In our application \(k=K\) is algebraically closed, so
that extension is trivial. The affinoid interpretation identifies
\(\mathrm{Sp}(A)\) with \(\mathcal M(A)\), bounded algebra maps with
analytic maps, and an affinoid subdomain of the rigid algebraic
analytification with the affinoid domain of its Berkovich
analytification. This is the affinoid-domain compatibility in Berkovich
1993, §1.6, and is precisely the interpretation used by Thuillier in the
inspected passage. We do not invoke an unspecified equivalence of
arbitrary nonparacompact spaces. All translation here is on one affinoid
and an algebraic curve. It preserves the affinoid algebra, its morphism,
and its analytic-domain structure. The primary input for this proof is
the \textbf{explicit Berkovich-domain assertion in Thuillier's proof
just quoted}, used only over algebraically closed \(K\), with its
embedding premise independently checked in van der Put. No separately
unchecked general equivalence theorem is needed as a premise. The
interpretation does not make the entire affinoid ordinary open; the
ordinary-interior step is separately proved below.

The cached primary PDF has source URL
\url{https://www.numdam.org/item/AIF_1980__30_4_155_0.pdf}, size 860722
bytes, and SHA-256:

\begin{Shaded}
\begin{Highlighting}[]
\NormalTok{583b012282ec23f8da13b9e430a088e48953983643dde6f8a23faf5f2e1b3aca}
\end{Highlighting}
\end{Shaded}

These match its local bytes and cached provenance. I inspected the PDF
images of printed pages 155--156 and the proof through page 159, rather
than relying solely on the OCR or a reference label. The earlier failed
retrieval in the blocked campaign supplies no theorem; this cached PDF
now removes that source gap.

\subsubsection{S5. Wanner comparison and the included-test
issue}\label{s5.-wanner-comparison-and-the-included-test-issue}

V. Wanner, arXiv:1801.04713v3, 8 November 2018, Section 4 fixes a
\textbf{smooth proper algebraic} curve over an \textbf{algebraically
closed complete nontrivially valued} field. Theorem 4.7 states: for an
ordinary open \(W\subset C^{\mathrm{an}}\), a continuous real-valued
function that is psh is Thuillier-subharmonic. Corollary 4.8 states the
equivalence for continuous real-valued functions. In that source,
Definition 3.16 means by psh positivity of the current \(D[u]\), not
membership in an approximation-generated hierarchy.

Here only Theorem 4.7's forward implication is needed. Its compact
scalar tests from Definition 3.11 have presentations in algebraic
very-affine unit moments. Those regular units are analytic units, so
each such test is an ordinary CLD test on the same ordinary open, with
the same compact support. The differentials commute with this inclusion.
The integration pairing is the canonical CLD pairing used explicitly
before Wanner Definition 3.16: \(D[u](\psi)=\int_W uD\psi\), described
there as a current in the sense of CLD. On these unit moments the
calibration is the canonical one; algebraic tropical chart integration
is that construction restricted to the algebraic moments. Thus a zero
equation against every ordinary analytic test implies Wanner-current
positivity for both signs. We do not need, and do not claim to prove,
equality of the entire algebraic and analytic test sheaves. Canonical
integration and its calibration are unchanged by an analytic isomorphism
and are compatible with analytic domains (CLD 14.3.5, 15.7.1--2 and
§7.7). In our comparison application the curves are smooth and reduced;
there is no multiplicity normalization or probability normalization to
change.

The cache is the stated arXiv v3 PDF; its bytes are not identified with
the journal PDF. The published citation is Math. Z. 293 (2019),
443--474, DOI 10.1007/s00209-018-2205-z. No Wanner theorem on arbitrary
analytic curves or on nonclosed fields is assumed.

\subsubsection{S6. Already-recognized harmonicity, not current
pullback}\label{s6.-already-recognized-harmonicity-not-current-pullback}

GR is W. Gubler and J. Rabinoff, arXiv:2506.13548v2, 17 September 2025.
The needed exact statements are:

\begin{itemize}
\tightlist
\item
  3.4.1--2: on a smooth connected strict analytic curve over a
  nontrivially valued field, a finite upper semicontinuous function is
  GR-subharmonic if and only if it is Thuillier-subharmonic. The
  identically \(-\infty\) alternative in 3.4.1 is irrelevant here.
\item
  3.12 and Definition 3.13: on a smooth strict curve, harmonicity agrees
  with both signs being subharmonic; on an arbitrary analytic curve,
  harmonicity is defined by that two-sign property. The proof in 3.12
  uses the strict affinoid Dirichlet solution (Thuillier Lemma 2.3.9)
  and comparison for both signs.
\item
  Proposition 3.14(1,6,7): harmonicity is a sheaf; it is preserved and
  detected by every analytic field extension; for any morphism of
  analytic curves, harmonicity pulls back, and the converse holds when
  that morphism is finite surjective. These statements have no
  reducedness or separability restriction. They concern functions
  \textbf{already} harmonic, not arbitrary continuous zero currents.
\item
  Lemma 3.9(1--2): the finite normalization of a reduced strict analytic
  curve is rig-smooth if the field is perfect, and is an isomorphism
  outside a discrete rigid set. We use the smoothness assertion only
  when \(K\) is algebraically closed and hence perfect; finiteness and
  boundarylessness make the normalization smooth in the good-space sense
  of S2.
\item
  4.1.2: on an analytic curve, classical psh equals GR-subharmonicity.
  Its extension-field and curve-map implication uses Proposition
  3.5(7,8), not a scalar-current assertion. Consequently a GR-harmonic
  curve function is classically pluriharmonic.
\item
  Proposition 3.15: on \textbf{any strict analytic curve} over a
  nontrivially valued field, a continuous function is harmonic if and
  only if it is real PL and every residual class
  \(L_u(x)\in\operatorname{Pic}(\operatorname{red}(Z,x))_{\mathbb R}\)
  is numerically trivial. This is precisely GJR harmonicity. GJR,
  arXiv:2111.05741v3, Definition 7.2 and Remark 7.3(1), define
  harmonicity by both-sign semipositivity and equate it to this
  numerical triviality for an R-PL function. Nonreduced curves are
  permitted in GR Proposition 3.15; its proof explicitly passes through
  reduction.
\end{itemize}

These sources justify the universal curve conclusions only after
recognition. In Theorem C the pullback is a pullback of a recognized
harmonic function; no arbitrary finite-map current pullback is imported.

\subsection{Smooth curve recognition on actual ordinary
neighbourhoods}\label{smooth-curve-recognition-on-actual-ordinary-neighbourhoods}

\textbf{Lemma 1 (ordinary local algebraic model).} Under the field
hypotheses of Theorem A, every point \(z\) of a good separated smooth
strict boundaryless pure curve \(Z\) has an ordinary neighbourhood
analytically isomorphic to an ordinary open of a smooth proper algebraic
curve's analytification over the same \(K\).

\textbf{Proof.} Choose a strict affinoid neighbourhood \(Q_0\) of the
arbitrary analytic point \(z\) inside \(Z\), using S2. It is normal and
pure one-dimensional by geometric regularity and domain invariance. An
affinoid algebra has finitely many connected components, represented by
idempotents. Select the component \(Q\) containing \(z\). Since it is
open and closed in \(Q_0\), it is still a strict affinoid neighbourhood
of \(z\) in \(Z\). Its algebra is a normal domain of dimension one: in a
noetherian normal ring distinct minimal primes cannot meet, because
every localization is a domain. The finitely many minimal components
therefore give an idempotent decomposition; connectedness leaves one.

S4 now supplies an isomorphism \(j:Q\simeq Q'\), where \(Q'\) is an
affinoid domain in \(C^{\mathrm{an}}\) for a nonsingular complete
algebraic curve \(C\) over \(K\). If necessary take its connected
component containing \(Q'\); \(C\) is a smooth proper algebraic curve.
This is a domain embedding, so an additional step is required.

Apply S3 to \(Q\hookrightarrow Z\). Since \(Z\) is separated and
boundaryless, \[
      \operatorname{int}_Z(Q)=\operatorname{Int}(Q)=Q\setminus\partial Q.
\] Apply S3 again to \(Q'\hookrightarrow C^{\mathrm{an}}\). The latter
is separated and boundaryless, so \[
      \operatorname{int}_{C^{\mathrm{an}}}(Q')
             =\operatorname{Int}(Q')=Q'\setminus\partial Q'.
\] Absolute analytic interior is invariant under the isomorphism \(j\).
Therefore \(j\) identifies these two \textbf{ordinary opens}. The first
contains \(z\) because \(Q\) was a neighbourhood. On both sides the open
analytic structure is the restriction of the domain structure, so the
induced map is an analytic isomorphism of ordinary neighbourhoods. This
proof treats every analytic point, including types three and four; it
uses neither a rigid-point-density cover nor an admissible cover in
place of an ordinary one. \(\square\)

\textbf{Lemma 2 (smooth recognition).} Under these hypotheses, a
continuous \(u\) with \(D[u]=0\) on \(Z\) is Thuillier harmonic and
GR-harmonic.

\textbf{Proof.} At any \(z\), use Lemma 1 to identify an ordinary
neighbourhood \(O\) with \(W\subset C^{\mathrm{an}}\). Restriction to
\(O\) is a legitimate restriction of a current. For a compact scalar
test \(\psi\) on \(W\), its pullback is a compact scalar test on \(O\)
and can be extended by zero to \(Z\); its support is compactly contained
in the open. Canonical integration and \(D\) are invariant under the
analytic isomorphism. Thus the transported continuous function
\(v=u\circ j^{-1}\) satisfies \(D[v]=0\) on \(W\) against every ordinary
test. In particular both signs satisfy Wanner's positive-current
hypothesis by S5. All hypotheses are now met: same algebraically closed
complete nontrivially valued field, smooth proper algebraic curve,
actual ordinary open, finite continuous function. Theorem 4.7 makes
\(v\) and \(-v\) Thuillier-subharmonic.

The two-sign criterion S6 makes \(v\) harmonic. Pull back the already
recognized function by the open analytic isomorphism. GR 3.4.2 and 3.12
identify this with Thuillier harmonicity on \(O\). The ordinary opens
\(O\) cover \(Z\) and harmonicity is a sheaf, so \(u\) is harmonic
throughout \(Z\). No current has been pulled through a scalar extension
or a nonisomorphic finite map. \(\square\)

\subsection{Singular and nonreduced intrinsic
recognition}\label{singular-and-nonreduced-intrinsic-recognition}

\textbf{Proof of Theorem A.} First suppose \(Z\) reduced. Its finite
normalization \(\nu:\widetilde Z\to Z\) is a good separated strict pure
curve. It is boundaryless since the finite map is boundaryless and \(Z\)
is boundaryless. The algebraically closed field is perfect, so S6 makes
it rig-smooth; S2 then makes it smooth. The normalization-transfer
theorem of n100 gives the actual equation \[
                         D[u\circ\nu]=0
\] on this normalization. Its hypotheses match exactly: reduced strict
boundaryless pure curve, complete nontrivially valued field, and a
continuous real multiplier. Lemma 2 recognizes \(u\circ\nu\) as
harmonic. GR Proposition 3.14(7), for the finite surjective curve
morphism \(\nu\), descends this \textbf{recognized} harmonicity to
\(u\). We have not deduced the upstairs equation by inversion of a
direct image.

For arbitrary \(Z\), the closed immersion
\(i:Z_{\mathrm{red}}\hookrightarrow Z\) is finite, surjective on
carriers, and its source is good separated strict boundaryless and pure
one-dimensional. Apply n100's closed-curve restriction with ambient
dimension \(n=1\) and \(C=Z_{\mathrm{red}}\). It yields
\(D[u\circ i]=0\) with the reduced source's canonical integration. This
explicitly respects the original equation's possibly unequal generic
nilpotent lengths; they have not been divided out. The reduced case
proves \(u\circ i\) harmonic. GR Proposition 3.14(7) applies to this
finite surjective morphism as well and descends harmonicity to \(u\) on
the original nonreduced curve.

Finally apply GR 4.1.2 to each sign. This gives classical
pluriharmonicity with the full extension-field quantifiers on the
\textbf{curve}. GR Proposition 3.15 gives local real-PL regularity and
numerical triviality of all residual classes, which is GJR harmonicity.
All constructions apply after ordinary-open restriction. This proves
Theorem A. No ordinary CLD smooth harmonic regularity or ordinary smooth
approximation is asserted. \(\square\)

\subsection{The affinoid branch lemma, including Banach
factorization}\label{the-affinoid-branch-lemma-including-banach-factorization}

\textbf{Lemma 3.} Let \(K\) be complete nontrivially valued,
\(\mathcal M(B)\) a connected normal strict affinoid pure curve, and
\(f:\mathcal M(B)\to\mathcal M(A)\) an analytic map to a strict
affinoid. Either its bounded algebra map factors through a finite field
\(E/K\) and a rigid point, or \(\dim_y f=0\) at every analytic source
point.

\textbf{Proof.} As in Lemma 1, connected normality makes \(B\) an
integral domain, of Krull dimension one. Fix an actual bounded Tate
presentation \(K\langle T_1,\ldots,T_N\rangle\twoheadrightarrow A\) and
write \(b_j=f^*(T_j)\in B\).

If every \(b_j\) is algebraic over \(K\), then
\(E=K[b_1,\ldots,b_N]\subset B\) is finite-dimensional over \(K\) and a
domain, hence a finite field. Every point of \(\mathcal M(B)\) restricts
to the unique extending absolute value on \(E\). Thus
\(\rho_B|_E=|\cdot|_E\). Since \(B\) is reduced its Banach norm is
equivalent to \(\rho_B\); \(E\) is therefore a closed complete
finite-dimensional subspace of \(B\) with its usual finite-field
topology. For every Tate series, polynomial partial sums evaluated at
the \(b_j\) converge in \(B\) under the bounded map and belong to \(E\).
Closedness puts the limit in \(E\). Consequently the \textbf{entire}
bounded map, not just its polynomial subalgebra, factors through \(E\).
Boundedness to \(E\) follows from \(|e|_E=\rho_B(e)\leq C\|e\|_B\) and
the bound of the original map; descent to \(A\) uses its quotient norm.
Its image contains every \(b_j\), so \(A\to E\) is surjective. The
inclusion \(E\to B\) is bounded by the same norm equivalence. Hence
\(f\) factors analytically through the singleton \(\mathcal M(E)\),
mapping to a rigid target point. No separability is required for
\(E/K\).

Otherwise some \(b=b_j\) is transcendental over \(K\). Let \(a\) be any
rigid source point. The element \(b(a)\in\mathcal H(a)\) is algebraic
over \(K\); let \(P_a\in K[t]\) be its monic minimal polynomial. The
element \(P_a(b)\in B\) is nonzero by transcendence, vanishes at \(a\),
and is not a unit. A quotient of a noetherian one-dimensional domain by
a nonzero nonunit has dimension zero: every prime containing it is
nonzero and therefore maximal. S2's strict Nullstellensatz and the
artinian argument make \(B/(P_a(b))\) finite over \(K\), retaining its
possible nilpotents.

Put \(x=f(a)\). The isometric embedding
\(\mathcal H(x)\hookrightarrow\mathcal H(a)\) makes \(x\) rigid, with
finite residue field over \(K\). In the actual analytic fibre
\(\mathcal M(B)\times_{\mathcal M(A)}\mathcal M(\mathcal H(x))\), the
coordinate \(b\) equals the evaluated target coordinate. Its polynomial
\(P_a\) is zero there, since it is zero after the injective
residue-field map into \(\mathcal H(a)\). The universal property of a
closed zero subspace therefore factors the fibre as a closed subspace of
\[
                 \mathcal M(B/(P_a(b)))_{\mathcal H(x)}.
\] This last space is finite over \(\mathcal H(x)\) and
zero-dimensional. Equivalently, the completed fibre algebra is the
quotient of the scalar-extended finite zero algebra by the target-value
relations. Thus \(\dim_a f=0\). This argument compares fibres over their
actual field and applies also when that field extension is inseparable.

The locus \(F=\{y:\dim_y f\geq1\}\) is Zariski closed by S3. On a strict
affinoid any nonempty such locus is the zero set of a proper ideal and
has a rigid point by S2. Such a point would contradict the preceding
calculation. Hence \(F\) is empty and \(\dim_y f=0\) for \textbf{every}
analytic \(y\). This uses Zariski upper semicontinuity, not an extension
by ordinary-topological density. \(\square\)

\subsection{Proof of the local image
theorem}\label{proof-of-the-local-image-theorem}

\textbf{Proof of Proposition B.} Fix \(y\in Y\) and \(x=f(y)\). Choose a
strict affinoid neighbourhood \(U_1\) of \(x\). Inside the ordinary open
\(f^{-1}(\operatorname{int}_X U_1)\) choose a strict affinoid
neighbourhood of \(y\) and then its connected component \(Q\) containing
\(y\). Smoothness makes its algebra normal and domain invariance makes
it pure one-dimensional, so Lemma 3 applies.

In its constant case, the bounded factorization through \(E\) restricts
to an ordinary neighbourhood of \(y\) and proves case 1. In the other
case the local relative dimension of the original \(f\) at \(y\) is
zero: the restriction fibre is an analytic domain neighbourhood in the
original fibre, and domain restriction preserves local dimension. Since
the target \(X\) is separated, S3 gives \[
                    \partial Y
                     =\partial(Y/X)\cup f^{-1}(\partial X).
\] The source is boundaryless, so \(y\) is in the relative interior of
\(f\). The finite-at-a-point criterion supplies analytic neighbourhoods
\(V'\to U'\) on which \(f\) is finite. This step establishes no
compactness of the original entire fibre and uses no such compactness.

Take a strict affinoid neighbourhood \(U\) of \(x\) inside an ordinary
neighbourhood contained in \(U'\). Its inverse image \(W\) in \(V'\) is
affinoid and finite over \(U\). It contains an ordinary neighbourhood of
\(y\) in \(Y\): intersect an ordinary neighbourhood contained in \(V'\)
with \(f^{-1}(\operatorname{int}_X U)\). Select the connected component
\(V\) of \(W\) containing \(y\). It is open and closed in the affinoid,
still a neighbourhood of \(y\), and still finite over \(U\). As a source
domain it is regular normal and pure one-dimensional, hence its algebra
\(B\) is a domain.

For completeness, the finite affinoid over the strict
\(A=\mathcal O(U)\) is strict here. Choose finitely many module
generators of its reduced algebra \(B\). Multiply each by a suitable
nonzero scalar from \(K\) so its spectral radius is at most one. Their
powers are bounded by the equivalence of reduced spectral and Banach
norms. They and the target Tate generators give a bounded surjection
from a strict Tate algebra onto \(B\). Its quotient topology is
equivalent to the finite-module topology: expressions as sums of the
chosen module generators bound the quotient norm by a fixed multiple of
the finite-module quotient norm, and the bounded surjection gives the
reverse bound. S3 identifies this finite-module topology with the
existing Banach topology. Thus \(V\) is strict. The same argument
applies to the finite reduced normalization algebra used below.

Let \(I=\ker(A\to B)\) and \(A_0=A/I\). It is a strict affinoid algebra
that is an integral domain, with an injective bounded map to \(B\), and \(B\) is finite
integral over \(A_0\). The closed subspace
\(C=\mathcal M(A_0)\hookrightarrow U\) is reduced. Injective integral
extensions preserve Krull dimension, by going-up and incomparability, so
\(\dim A_0=\dim B=1\). The prime kernel and S2 make \(C\) irreducible
and pure one-dimensional.

To check surjectivity at \textbf{all analytic points}, let \(c\in C\)
have kernel \(\mathfrak p\subset A_0\). Lying-over for the integral
injection gives a prime of \(B\) over \(\mathfrak p\). Thus
\(B\otimes_{A_0}\kappa(\mathfrak p)\) is a nonzero finite-dimensional
algebra, and extending its field to \(\mathcal H(c)\) leaves it nonzero.
Finite-module completion identifies this algebra with the finite
analytic fibre algebra \(B\widehat\otimes_{A_0}\mathcal H(c)\) (S3);
finite-dimensional vector spaces over a complete field are already
complete. A maximal ideal gives a finite residue field extension of
\(\mathcal H(c)\) with its extending absolute value, hence an analytic
point of the fibre. This proves that \(V\to C\) is finite surjective. It
is not an inference merely from surjectivity of the two separate maps to
algebra spectra.

Let \(\overline A_0\) be the finite integral closure of \(A_0\) in its
fraction field. It exists by S2. The injection into the domain \(B\)
embeds \(\operatorname{Frac}(A_0)\) in \(\operatorname{Frac}(B)\). An
element of \(\overline A_0\) is integral over \(A_0\), hence integral
over \(B\), and belongs to \(B\) by normality. This gives the algebraic
lift \(\overline A_0\to B\). It is bounded: choose finitely many
\(A_0\)-module generators for \(\overline A_0\); boundedness of
\(A_0\to B\) and the norms of their images bound the module map from a
free presentation, and hence its quotient norm. S3 gives the finite
Banach topology on \(\overline A_0\). Further, the same \(A_0\)-module
generators of \(B\) generate \(B\) over \(\overline A_0\), so the lift
is finite as a ring map; the strict nontrivial finite-Banach assertion
of S3 makes it finite analytic. This is the asserted lift to the
normalization. It does not require a separable map, a geometrically
reduced image, or finite flatness. \(\square\)

\subsection{Intrinsic equations on local images and ground-field
probes}\label{intrinsic-equations-on-local-images-and-ground-field-probes}

\textbf{Proof of Theorem C.} At \(y\in Y\), apply Proposition B. In the
constant branch \(h\circ f\) is locally constant and therefore harmonic.
Otherwise set \(U^o=\operatorname{int}_X U\), an ordinary neighbourhood
of \(x=f(y)\), and \(C^o=C\cap U^o\). The open \(U^o\) is good separated
strict boundaryless and pure of dimension \(n\), and \(C^o\) is a closed
reduced pure curve inside it. Boundarylessness follows from the
boundaryless open inclusion and the finite closed immersion. Domain
invariance gives the stated dimensions. Restrict the ambient equation to
\(U^o\) and apply n100's closed-curve theorem, obtaining \[
                         D[h|_{C^o}]=0.
\] Its normalization-transfer theorem then gives
\(D[h|_{C^o}\circ\nu]=0\) on the finite normalization of \(C^o\). All
these equations are intrinsic, on the actual ordinary opens, with
canonical integration. No test meets the affinoid target boundary.
Normalization restricts to \(C^o\) because integral closure is local.

This normalization is a good separated smooth strict boundaryless pure
curve over the perfect field \(K\). Lemma 2, or the smooth case of
Theorem A, recognizes the lifted function as GR-harmonic. The lift in
Proposition B maps the ordinary neighbourhood \[
         \operatorname{int}_Y(V)\cap f^{-1}(U^o)
\] of \(y\) into that normalization. Pull back the \textbf{recognized
harmonic function} using GR Proposition 3.14(7). It gives harmonicity of
\(h\circ f\) near \(y\), irrespective of inseparability of the local
map. On the smooth strict source this is Thuillier harmonicity by S6.
Locality treats every analytic point of \(Y\), proving the theorem.

We never invert a finite-current direct image. We never replace a curve
image by its graph in an unproved product-current ascent. The finite
image, its ordinary interior, its normalization equation, and
recognition precede harmonic pullback. \(\square\)

\subsection{Comparison with the catalogue and the remaining
gates}\label{comparison-with-the-catalogue-and-the-remaining-gates}

The catalogue's continuous curve-recognition scope is an ordinary open
in the analytification of a smooth proper algebraic curve over an
algebraically closed field. Theorem A reaches \textbf{all good separated
strict boundaryless pure analytic curves over that field}, including
singular and nonreduced curves, with no global proper algebraic model.
The smooth step is local and reduces to the published comparison using
the checked affinoid embedding and the separately proved ordinary
interior identification. The singular and nilpotent step uses exactly
n100. This is a substantive recognition theorem in a new stated scope,
but is a special case of the original atom, not a full solution.

Proposition B establishes the proposed local-image route for smooth
strict boundaryless \textbf{ground-field} probes over arbitrary complete
nontrivially valued fields. In particular, a nonconstant branch has an
actual local analytic curve image even when global algebraic coordinate
relations need not exist. It claims no arbitrary finite-cover current
pullback. Theorem C combines this image result and recognition over
algebraically closed \(K\). Its exact extra content is harmonicity of
every such \(K\)-probe into a pure ambient space.

For general \(K\), the normalization of a reduced curve need not be
rig-smooth before treating geometrically inseparable structure. Even
over a perfect nonclosed field, normal rig-smooth curves still require
current recognition there. Neither van der Put's algebraically closed
field hypothesis nor Wanner's hypothesis can be dropped. The reviewed
purely inseparable ascent packet mentioned by the operator concerns
characteristic \(p\) and the completion of a specified purely
inseparable algebraic extension only. No result from that packet is a
premise here. Its committed exact-scope PASS is acknowledged without
transferring it to changed bytes or to separable/transcendental ascent.
Completed perfection is perfect but may not be algebraically closed.

The \textbf{scalar-current ascent gate} remains distinct. GR Definition
4.1 requires all analytic \(L/K\) and curve maps into \(X_L\). Theorem C
supplies none of the missing equations on these actual scalar-extended
spaces. Algebraic closure of \(K\) removes finite algebraic extensions,
not transcendental analytic extensions. Theorem A's universal classical
conclusion is legitimate because harmonicity has already been recognized
on a curve, where GR 4.1.2 and 3.14 provide the full quantifiers; it
cannot be applied to turn a fieldwise ambient test property into that
same universal conclusion.

The \textbf{two-sign ordinary regularization gate} also remains
separate. This proof does not construct ordinary CLD-smooth harmonic
approximants or decreasing ordinary smooth psh approximants on arbitrary
singular curves or arbitrary ambient germs. I read the controlling book
chapter and the complete frozen P06.3 statement and proof. Their
endpoint assumes algebraically closed \(K\) and an ordinary open of a
smooth proper algebraic curve; its recognition implication is Wanner
Corollary 4.8. Its other direction constructs psh approximants using
finite tree skeleta, positive rational strictification, residual-divisor
movement into fresh directions, convex caps, and a decreasing tail
shift. The approximants can carry positive curvature, so this is not an
ordinary smooth harmonic-approximation theorem. That packet is not a
supporting theorem in the present proof. The present ordinary local
models could support a separate smooth-curve regularization
globalization after its own source audit; no new regularization theorem
is claimed here.

No counterexample satisfying the original strict ambient hypotheses is
constructed. No conclusion is drawn for arbitrary nonproduct
higher-dimensional germs, nonclosed fields, mixed-dimensional ambient
current conventions, or the non-strict invisible-boundary product. The
status is \texttt{partial} and \texttt{candidate=false}.

\subsection{Repair, self-audit, and
reproducibility}\label{repair-self-audit-and-reproducibility}

The earlier submission's native support entries contained citation
prose. This repaired submission uses only the supplied ID n100 or
earlier new local IDs n1, n2, \ldots{} in support arrays. Citations,
paths, hashes, and commands appear in note text, not support. The prior
invalid findings, argument, audit, and research notes are preserved
locally as diagnostic data; no old campaign, database, settings, or
proof file was modified.

The mathematical change is substantive: the previous draft's smooth
recognition proof used semistable factorization, skeletal integration,
and vertex slope detection. The new proof removes those load-bearing
steps and replaces them with the now-inspected primary affinoid
embedding theorem, two explicit boundary equalities, included compact
tests, and the published comparison at exactly its stated scope. The
local image proof is rechecked here, including constant-branch Banach
limits, actual residue-field fibres, all-point Zariski reasoning, strict
finite affinoids, spectral surjectivity, and bounded finite
normalization lifts. No preceding self-audit or verdict substitutes for
these arguments.

\texttt{self-audit.md} records applicability and remaining obligations.
\texttt{research-notes.md} records retained and discarded routes. The
standalone paper is \texttt{P06.2-partial.tex}, with
\texttt{P06.2-partial.pdf}. The requested five-field record is
\texttt{status.json}. The programs, source copies, commands, outputs,
environment versions, and manifests are saved here.

Rerun from the preserved worker workspace
\path{/workspace/research/campaigns/round-04/P06.2-curve-recognition-submission-repair/artifacts/codex-bcowsD}:

\begin{Shaded}
\begin{Highlighting}[]
\ExtensionTok{python}\NormalTok{ prepare.py}
\ExtensionTok{python}\NormalTok{ build{-}and{-}check.py}
\end{Highlighting}
\end{Shaded}

The second command generates the standalone TeX with installed pandoc,
runs installed pdfLaTeX twice, records outputs, checks the native note
dependency/schema rules, and writes artifact hashes. Compilation and
finite source/schema diagnostics establish reproducibility and
typesetting only; they are not mathematical verification. No
installation, extra model call, subagent, heavy computation, external
message, or environment/authentication modification is used. No
resource-heavy computation is needed for these proofs. The standalone
publication source is also compiled twice with installed pdfLaTeX.
The native editor connection did not return, so built-in compilation
remains unverified; the source and independently compiled PDF are
preserved.

\subsection{Bibliography}\label{bibliography}

\begin{enumerate}
\def\labelenumi{\arabic{enumi}.}
\tightlist
\item
  M. van der Put, \emph{The class group of a one-dimensional affinoid
  space}, Ann. Inst. Fourier 30(4) (1980), 155--164, Theorem 1.1; field
  and affinoid conventions on p.~155; proof on pp.~156--159.
\item
  V. Wanner, \emph{Comparison of two notions of subharmonicity on
  non-archimedean curves}, arXiv:1801.04713v3; Math. Z. 293 (2019),
  443--474. Definitions 3.11, 3.16, Theorem 4.7, Corollary 4.8.
\item
  A. Thuillier, \emph{Théorie du potentiel sur les courbes en géométrie
  analytique non archimédienne. Applications à la théorie d'Arakelov},
  thesis, 2005, HAL tel-00010990v1. Theorem 2.3.8 and its actual
  affinoid-embedding assertion; Lemma 2.3.9; Definition 3.1.5.
\item
  A. Chambert-Loir and A. Ducros, \emph{Formes différentielles réelles
  et courants sur les espaces de Berkovich}, arXiv:1204.6277v2, 25 July
  2025. §§7.7, 11.1--4, 11.7.7, 14.3.5, 15.7.1--2.
\item
  W. Gubler and J. Rabinoff, \emph{Classically psh and pluriharmonic
  functions on Berkovich spaces}, arXiv:2506.13548v2, 17 September 2025.
  §§3.4.1--2, 3.9, 3.12--15, 4.1.2.
\item
  W. Gubler, P. Jell, and J. Rabinoff, \emph{Forms on Berkovich spaces
  based on harmonic tropicalizations}, arXiv:2111.05741v3, 7 March 2025.
  Definition 7.2 and Remark 7.3(1). This inspected snapshot labels
  itself final published version; no later numbering is silently
  substituted.
\item
  B. Conrad and M. Temkin, \emph{Descent for non-Archimedean analytic
  spaces}, arXiv:1912.06230v2, Lemma 2.9, inspected author-hosted PDF.
\item
  V. Berkovich, \emph{Étale cohomology for non-Archimedean analytic
  spaces}, Publ. Math. IHÉS 78 (1993), §1.6, affinoid-domain
  compatibility with rigid geometry, as used in the inspected Thuillier
  passage.
\item
  S. Bosch, U. Güntzer, and R. Remmert, \emph{Non-Archimedean Analysis},
  Springer, 1984, 3.2.4/2 and 6.1.2/1; the finite Banach formulation is
  read in CLD 11.4.4 and its footnote.
\item
  A. Ducros, \emph{Variation de la dimension relative en géométrie
  analytique p-adique}, Compositio Math. 143 (2007), 1511--1532, Theorem
  4.9, in the exact formulation read in CLD 11.4.5.
\item
  Supplied note n100, \emph{Intrinsic zero-Hessian restriction to closed
  analytic curves}, frozen complete argument and exact-scope review;
  provenance and exact hash stated in S1. Used only at its declared
  scope.
\end{enumerate}

\end{document}
